\documentclass[10pt,a4paper]{amsart}
\usepackage[margin=19mm]{geometry}
\usepackage{amsmath,amssymb,mathtools}
\usepackage[scr=rsfs,cal=euler]{mathalfa}
\usepackage{xfrac}
\usepackage[unicode,hidelinks,bookmarks=false]{hyperref}
\newcommand{\ie}{i.e.\ }
\newcommand{\cf}{cf.\ }
\newcommand{\resp}{respectively\ }
\newcommand{\Subsection}{\subsection}
\providecommand{\nobreakdash}{\nobreak\hbox{-}}
\setlength{\emergencystretch}{2em}
\multlinegap0pt
\usepackage{bm}
\usepackage{bbm}
\usepackage{dsfont}
\usepackage{xspace}
\usepackage{cancel}
\usepackage{mathtools}
\usepackage{amsthm}
\makeatletter
\@ifundefined{lemma}{\newtheorem{lemma}{Lemma}}{}
\makeatother
\DeclareMathOperator*{\argmin}{arg\,min}
\newcommand{\bbE}{\mathbb{E}}
\newcommand{\bbR}{\mathbb{R}}
\newcommand{\bbZ}{\mathbb{Z}}
\newcommand{\calI}{\mathcal{I}}
\newcommand{\calN}{\mathcal{N}}
\newcommand{\calX}{\mathcal{X}}
\newcommand{\dmin}{\mathsf{d}_{\min}}
\newcommand{\intR}{\int_{\bbR}}
\newcommand{\snr}{\mathsf{snr}}
\renewcommand{\d}{\,\mathrm{d}}
\title[High signal-to-noise ratio asymptotics of entropy-constraine]{High signal-to-noise ratio asymptotics of entropy-constrained Gaussian channel capacity}
\author{Adway Girish, Shlomo Shamai, Emre Telatar}
\date{Source: arXiv:2601.09864v1 (CC BY 4.0)}
\begin{document}
\maketitle

\noindent\emph{Source and licence.} High signal-to-noise ratio asymptotics of entropy-constrained Gaussian channel capacity. Version arXiv:2601.09864v1 (CC BY 4.0), distributed under the Creative Commons Attribution 4.0 International licence, as recorded in the arXiv licence metadata for this exact version and shown on the abstract page https://arxiv.org/abs/2601.09864v1. Full rights notice: licenses/arxiv-2601.09864-CC-BY-4.0.txt.\par\smallskip

\noindent\textit{Editorial scope.} Counted unit: the Lemma whose source label is \texttt{lem: cond\_ent} in the article (source0.tex lines 269--273, source label \texttt{lem: cond\_ent}), delivered together with the article's own complete original proof of it (source0.tex lines 274--400, one \texttt{proof} environment of 127 lines). No mathematics is written, repaired or re-derived: statement and proof are the article's own text line for line. Required context retained uncounted: none. Every definition, display and label of those lines is retained unchanged. Omitted from this package and not redistributed: 1--268, 401--736. Nothing inside a retained excerpt is omitted or altered: every retained body is delivered exactly as the article's lines are written, so no omission receipt of any kind accompanies this package. Citations: the retained excerpts carry no citation command at all, so no bibliography entry of the article is transcribed and no reference is redistributed. Reference adaptation: none was needed and none was made; every reference inside the delivered unit resolves inside it, so no unresolved reference or citation is produced. Printed slips kept literal: no printed word, symbol, hypothesis or number of the counted statement or of its proof was altered. Layout and class: the article's own class is \texttt{article} with packages including geometry, inputenc, fontenc, hyperref, url, booktabs, amsfonts, amsmath, dsfont, amsthm, mathtools, amssymb, subcaption, pgfplots; the wrapper is the lane's pinned \texttt{amsart} editorial wrapper at 10pt on A4, which restates the theorem-like environments the retained text needs and adds the article's own macro definitions that the retained text needs (\texttt{\textbackslash argmin}, \texttt{\textbackslash bbE}, \texttt{\textbackslash bbR}, \texttt{\textbackslash bbZ}, \texttt{\textbackslash calI}, \texttt{\textbackslash calN}, \texttt{\textbackslash calX}, \texttt{\textbackslash d}, \texttt{\textbackslash dmin}, \texttt{\textbackslash intR}, \texttt{\textbackslash snr}), with extra wrapper packages (bm, bbm, dsfont, xspace, cancel, mathtools). The delivered numbering is the unit's own. Assets: no retained span contains a figure environment, a graphics-inclusion command, a PDF-page inclusion, a table or a TikZ picture; no image file is copied into this package, the package declares no asset dependency and no image file is redistributed with it. Licence: version arXiv:2601.09864v1 of this article is distributed under the Creative Commons Attribution 4.0 International licence, as recorded in the arXiv licence metadata for this exact version and shown on the abstract page https://arxiv.org/abs/2601.09864v1. A full rights notice is delivered at licenses/arxiv-2601.09864-CC-BY-4.0.txt. Characters: every retained excerpt is pure ASCII, so the delivered engine, XeLaTeX with its default Latin Modern text font, reports no missing character.
	\begin{lemma}[Exponent of conditional entropy] For any discrete $X$ with a finite entropy, and $Y = X + Z$ with $Z \sim \calN(0,1/\snr)$ independent of $X$, we have
		\begin{equation*}
			\lim_{\snr \to \infty} - \frac{1}{\snr} \log H(X \mid Y) = \frac{\dmin(X)^2}{8}.
		\end{equation*} \label{lem: cond_ent}
	\end{lemma}

	\begin{proof}[Proof]
		Suppose $X$ is discrete and has entropy $H(X) < \infty$. We will show the above result by deriving upper and lower bounds to $H(X \mid Y)$ that have the same exponential rate in $\snr \eqqcolon \frac{1}{\sigma^2}$. 
		We also use the following notation: let $\varphi_{\sigma}(x)$ denote the density of a Gaussian random variable with mean zero and variance $\sigma^2$, i.e., $\varphi_{\sigma}(x) \coloneqq \frac{1}{\sqrt{2\pi\sigma^2}} \exp(-\frac{x^2}{2\sigma^2})$, and let $\varphi \coloneqq \varphi_1$ denote the standard Gaussian density.
		
		\paragraph{(a) Lower bound to $H(X \mid Y)$. } 
		Fix $\epsilon > 0$. 
		Let $x_0$ and $x_1$ be such that $x_0 < x_1$ and $\Delta = x_1 - x_0 < \dmin + \epsilon$. Note that $x_0,x_1,\Delta$ depend on $\epsilon$ but we do not write this dependence explicitly. Since $H(X \mid Y)$ is the average of $\log\frac1{p_{X|Y}}(x|y)$ (which is always non-negative) over the joint distribution $(X,Y)$, we can lower bound $H(X \mid Y)$ as
		\begin{align*}
			H(X \mid Y) &= \sum_{x \in \calX}  \intR \log \frac{\sum_{x' \in \calX} p_{x'} \varphi_{\sigma}(y-x')}{p_{x} \varphi_{\sigma}(y - x)}\ p_{x}\varphi_{\sigma}(y-x) \d y\\
			&\geq  \intR \log \left(1+\frac{ p_{x_1} \varphi_{\sigma}(y-x_1)}{p_{x_0} \varphi_{\sigma}(y - x_0)}\right) p_{x_0}\varphi_{\sigma}(y-x_0) \d y,
		\end{align*}
		by considering only the $x_0$ term in the outer sum and only the $x_0$ and $x_1$ terms in the sum inside $\log$. 
		We now further lower bound this integral by picking an arbitrary $\delta > 0$ and restricting the domain of integration to the interval $(\frac{x_0 + x_1}{2}, \frac{x_0 + x_1}{2}+\delta)$, i.e.,
		\begin{align*}
			H(X \mid Y) &\geq  \int_{\frac{x_0 + x_1}{2}}^{\frac{x_0 + x_1}{2}+\delta} \log \left(1+\frac{ p_{x_1} \varphi_{\sigma}(y-x_1)}{p_{x_0} \varphi_{\sigma}(y - x_0)}\right) p_{x_0}\varphi_{\sigma}(y-x_0) \d y.
		\end{align*}
		Consider the expression $$\frac{ p_{x_1} \varphi_{\sigma}(y-x_1)}{p_{x_0} \varphi_{\sigma}(y - x_0)} = \frac{ p_{x_1} \exp(-\frac{(y-x_1)^2}{2\sigma^2})}{p_{x_0} \exp(-\frac{(y-x_0)^2}{2\sigma^2})} = \frac{p_{x_1}}{p_{x_0}}\exp\left(\frac{1}{\sigma^2}(x_1-x_0)(y - \tfrac{x_0 + x_1}{2})\right).$$
		This is increasing in $y$ as $x_1 > x_0$, so we can obtain a lower bound to the above expression by substituting $y = \frac{x_0 + x_1}{2}$, which makes the argument of the exponent vanish, giving 
		\begin{align*}
			H(X \mid Y) &\geq p_{x_0}\log \left(1+\tfrac{p_{x_1}}{p_{x_0}} \right) \int_{\frac{x_0 + x_1}{2}}^{\frac{x_0 + x_1}{2}+\delta}  \varphi_{\sigma}(y-x_0) \d y \\
			&= p_{x_0}\log \left(1+\tfrac{p_{x_1}}{p_{x_0}} \right) \int_{\frac{\Delta}{2\sigma}}^{\frac{\Delta}{2\sigma} + \frac{\delta}{\sigma}}  \varphi(y) \d y,
		\end{align*}
		with the second line obtained by making the substitution $\frac{y-x_0}{\sigma} \mapsto y$. 
		The integrand $\varphi(y)$ is decreasing for $y$ in the domain of integration (positive $y$), so we can lower bound the integrand by substituting $y = \frac{\Delta}{2\sigma} + \frac{\delta}{\sigma}$, i.e.
		\begin{align*}
			H(X \mid Y) &\geq p_{x_0}\log \left(1+\tfrac{p_{x_1}}{p_{x_0}} \right) \cdot \tfrac{1}{\sqrt{2\pi \sigma^2}} \exp\left(- \tfrac{(\Delta + 2\delta)^2}{8\sigma^2}\right) \cdot \tfrac{\delta}{\sigma},
		\end{align*}
		for every $\delta, \epsilon > 0$. Hence, the logarithm of the conditional entropy is lower bounded as
		\begin{align*}
			\log H(X \mid Y) \geq \log \left[\tfrac{1}{\sqrt{2\pi \sigma^2}}\tfrac{\delta}{\sigma} p_{x_0}\log \left(1+\tfrac{p_{x_1}}{p_{x_0}} \right) \right] - \tfrac{(\Delta + 2\delta)^2}{8\sigma^2},
		\end{align*}
		which implies that the exponential rate as $\sigma \to 0$ (i.e., $\snr \to \infty$) is 
		\begin{align*}
			\lim_{\sigma \to 0} \sigma^2 \log H(X \mid Y) &\geq \lim_{\sigma \to 0} \sigma^2\log \left[\tfrac{1}{\sqrt{2\pi \sigma^2}}\tfrac{\delta}{\sigma} p_{x_0}\log \left(1+\tfrac{p_{x_1}}{p_{x_0}} \right) \right] - \tfrac{(\Delta + 2\delta)^2}{8}\\
			&= -\tfrac{(\Delta + 2\delta)^2}{8},
		\end{align*}
		for every $\delta, \epsilon > 0$. Letting $\delta \to 0$, we have 
		$$\lim_{\sigma \to 0} \sigma^2 \log H(X \mid Y) \geq -\tfrac{\Delta^2}{8} > -\tfrac{(\dmin + \epsilon)^2}{8},$$
		as $\Delta < \dmin + \epsilon$, by our choice of $x_0, x_1$. Now letting $\epsilon \to 0$, we have the desired result, i.e., $\lim_{\sigma \to 0} \sigma^2 \log H(X \mid Y) \geq -\tfrac{\dmin^2}{8}$. (Note that we do not require $\dmin > 0$ anywhere, the proof also works for $\dmin = 0$.)
		
		
		\paragraph{(b) Upper bound to $H(X \mid Y)$. }
		\sloppy
		
		If $\dmin = 0$, then we use the trivial upper bound $H(X \mid Y) \leq H(X)$, as this gives $\lim_{\sigma \to 0}  \sigma^2\log H(X \mid Y) \leq \lim_{\sigma \to 0} \sigma^2 \log H(X) = 0 = -\dmin^2/8$. 
		For the remainder of this upper bound, assume that $\dmin > 0$. 
		Hence, the support is either finite, or if it is infinite, there are no finite accumulation points. 
		
		In either case, let $\calX = (x_i)_{i \in \calI}$ be the set of atoms of $X$, with $\calI = \{1,\dots,n\}$ if $X$ is finite and $\calI = \bbZ$ if $X$ is infinite. We may, without loss of generality, assume that $x_i < x_{i+1}$ for all $i$, as there are no accumulation points. (This may not be possible if the atoms of $X$ have an accumulation point, consider the example $\calX = \{0\} \cup \{-1-1/n : n \in \bbZ\} \cup \{1 + 1/n : n \in \bbZ\}$.)
		
		
		
		
		By definition, we have $h(Y) = -\intR  f_Y(y) \log f_Y(y) \d y$, with $f_Y(y) = \sum_{x \in \calX} p_x  \frac{1}{\sqrt{2\pi \sigma^2}} \exp\left(-\frac{(y-x)^2}{2\sigma^2}\right) = \sum_{i \in \calI} p_i \frac{1}{\sqrt{2\pi \sigma^2}} \exp\left(-\frac{(y-x_i)^2}{2\sigma^2}\right)$ being the density of $Y$. This can be simplified as
		\begin{align*}
			h(Y) &= -\intR  f_Y(y) \log \left[\sum_{j \in \calI} p_j \tfrac{1}{\sqrt{2\pi \sigma^2}}\exp\left(-\frac{(y-x_j)^2}{2\sigma^2}\right)\right] \d y\\
			&= -\intR f_Y(y) \log \left[\sum_{j \in \calI} p_j \exp\left(-\frac{(y-x_j)^2}{2\sigma^2}\right)\right] \d y - \intR f_Y(y) \log\tfrac{1}{\sqrt{2\pi \sigma^2}}\d y\\
			&= \frac12\log(2\pi\sigma^2)-\intR \sum_{i \in \calI} p_i \tfrac{1}{\sqrt{2\pi \sigma^2}}\exp\left(-\frac{(y-x_i)^2}{2\sigma^2}\right) \log \left[\sum_{j \in \calI} p_j \exp\left(-\frac{(y-x_j)^2}{2\sigma^2}\right)\right] \d y,
		\end{align*}
		as $f_Y$ integrates to 1.  For each $y$, let $i^*(y) \in \calI$ be the index of the atom of $X$ that is closest to $y$, i.e., $i^*(y) = \argmin_i |y - x_i|$ (in case of ties, choose one arbitrarily --- there will only be countably many such $y$'s; a minimizing $i^*$ is also guaranteed to exist for all $y \in \bbR$ as $\dmin > 0$). Also let $d^*(y) = y - x_{i^*(y)}$ and $d_i(y) = y - x_i$. We write simply $i^*$, $d^*$ and $d_i$ when $y$ is clear from context. 
		
		As $\sum_{j \in \calI} p_j = 1$, we have that $\sum_{j \in \calI} p_j \exp\left(-\tfrac{d_j^2}{2\sigma^2}\right)  \leq 1$, and the integrand above is negative.
		Then, we can lower bound $h(Y)$ by considering only the term corresponding to $i = i^*(y)$ in the summation over $i$, to obtain
		\begin{align*}
			h(Y) &= \frac12\log(2\pi\sigma^2) - \intR p_{i^*} \tfrac{1}{\sqrt{2\pi \sigma^2}}\exp\left(-\tfrac{{d^*}^2}{2\sigma^2}\right) \log \left[\sum_{j \in \calI} p_j \exp\left(-\tfrac{d_j^2}{2\sigma^2}\right) \right]\d y.
		\end{align*}
		Note that $H(X \mid Y) = H(X) - I(X;Y) = H(X) - h(Y) + h(Y \mid X) = H(X) + \frac{1}{2}\log(2\pi e \sigma^2) - h(Y)$, as $h(Y \mid  X) = h(Z)$, which gives
		\begin{align}
			H(X \mid Y) &= H(X) + \frac12 + \intR p_{i^*} \tfrac{1}{\sqrt{2\pi \sigma^2}}\exp\left(-\tfrac{{d^*}^2}{2\sigma^2}\right) \log \left[\sum_{j \in \calI} p_j \exp\left(-\tfrac{d_j^2}{2\sigma^2}\right) \right]\d y . \label{eqn: h_x_y}
		\end{align}
		We re-write the integral by splitting $\log \left[\sum_{j \in \calI} p_j \exp\left(-\tfrac{d_j^2}{2\sigma^2}\right)\right]$ as $\log \left[p_{i^*} \exp\left(-\tfrac{{d^*}^2}{2\sigma^2}\right)\right] + \log \left[1 + \sum_{j\in \calI \setminus \{i^*\}} \frac{p_j}{p_{i^*}} \exp\left(\tfrac{{d^*}^2-d_j^2}{2\sigma^2}\right)\right]$. Hence, we have
		\begin{align}
			H(X \mid Y) &= H(X) + \frac12 + \intR p_{i^*} \tfrac{1}{\sqrt{2\pi \sigma^2}}\exp\left(-\tfrac{{d^*}^2}{2\sigma^2}\right) \log \left[p_{i^*} \exp\left(-\tfrac{{d^*}^2}{2\sigma^2}\right)\right]\d y  \nonumber \\
			&\qquad + \intR p_{i^*} \tfrac{1}{\sqrt{2\pi \sigma^2}}\exp\left(-\tfrac{{d^*}^2}{2\sigma^2}\right) \log \left[1 + \sum_{j\in \calI \setminus \{i^*\}} \frac{p_j}{p_{i^*}} \exp\left(\tfrac{{d^*}^2-d_j^2}{2\sigma^2}\right)\right] \d y. \label{eqn: h_x_y_ub}
		\end{align}
		
		To compute these integrals, define $\Delta_i = x_i - x_{i-1}$ for $i \in \bbZ$ when the support is infinite.  When $|\calX| = n$, define $\Delta_i = x_i - x_{i-1}$ for $i = 2,\dots,n$, $\Delta_0 = -\infty$ and $\Delta_n = \infty$.
		Also define $V_i$ to be the open interval $(x_i - \Delta_i/2, x_i + \Delta_{i+1}/2) = (\frac{x_i + x_{i-1}}{2}, \frac{x_{i+1} + x_{i}}{2})$ for $i \in \calI$. Observe that we have $i^*(y) = i$ for all $y$ in $V_i$, and that the $V_i$'s partition $\bbR$.
		We thus compute the first integral in \eqref{eqn: h_x_y_ub}, as
		\begin{align*}
			&\quad \ \intR p_{i^*} \tfrac{1}{\sqrt{2\pi \sigma^2}}\exp\left(-\tfrac{{d^*}^2}{2\sigma^2}\right) \log \left[p_{i^*} \exp\left(-\tfrac{{d^*}^2}{2\sigma^2}\right)\right]\d y\\
			&= \sum_{i \in \calI} \int_{V_i} p_{i} \tfrac{1}{\sqrt{2\pi \sigma^2}}\exp\left(-\tfrac{{d_i}^2}{2\sigma^2}\right) \log \left[p_{i} \exp\left(-\tfrac{{d_i}^2}{2\sigma^2}\right)\right]\d y\\
			&= \sum_{i \in \calI} \left\{ p_i\log p_i \int_{V_i}  \tfrac{1}{\sqrt{2\pi \sigma^2}}\exp\left(-\tfrac{(y-x_i)^2}{2\sigma^2}\right) \d y - p_i\int_{V_i} \frac{(y-x_i)^2}{2\sigma^2}\tfrac{1}{\sqrt{2\pi \sigma^2}}\exp\left(-\tfrac{(y-x_i)^2}{2\sigma^2}\right) \d y \right\}\\
			&= \sum_{i \in \calI} \left\{ p_i\log p_i \left[Q\left(\tfrac{-\Delta_i}{2\sigma}\right)-Q\left(\tfrac{\Delta_{i+1}}{2\sigma}\right)\right] - \frac{p_i}{2}\left[R\left(\tfrac{-\Delta_i}{2\sigma}\right)-R\left(\tfrac{\Delta_{i+1}}{2\sigma}\right)\right] \right\},
		\end{align*}
		where $Q$ is the standard $Q$-function and $R(z) = \int_z^{\infty} \frac1{\sqrt{2\pi}}u^2\exp(-\frac{u^2}{2}) \d u$.
		The $Q$-function satisfies the following useful properties: For $z > 0$, $Q(z)$ satisfies $\frac{z^2}{1+z^2} \frac{1}{\sqrt{2\pi}z}\exp(-\frac{z^2}{2}) \leq Q(z) \leq \frac{1}{\sqrt{2\pi}z}\exp(-\frac{z^2}{2})$. We also have, by definition, that $Q(-z) = 1-Q(z)$.
		The $Q$-function can be related to $R(z)$ by integrating the latter by parts to obtain $R(z) = Q(z) + \frac{z}{\sqrt{2\pi}}\exp(-\frac{z^2}{2})$. Hence, the above expression is equal to 
		\begin{align*}
			&\quad \ \sum_{i \in \calI} \left\{ \left(p_i\log p_i-\tfrac{p_i}{2}\right) \left[Q\left(\tfrac{-\Delta_i}{2\sigma}\right)-Q\left(\tfrac{\Delta_{i+1}}{2\sigma}\right)\right] - \tfrac{p_i}{2\sqrt{2\pi}}\left[\tfrac{-\Delta_i}{2\sigma}\exp\left(\tfrac{-\Delta_i^2}{8\sigma^2}\right)-\tfrac{\Delta_{i+1}}{2\sigma}\exp\left(\tfrac{-\Delta_{i+1}^2}{8\sigma^2}\right)\right] \right\} \\ 
			&= \sum_{i \in \calI} \left\{ \left(p_i\log p_i-\tfrac{p_i}{2}\right) \left[1-Q\left(\tfrac{\Delta_i}{2\sigma}\right)-Q\left(\tfrac{\Delta_{i+1}}{2\sigma}\right)\right] + \tfrac{p_i}{2\sqrt{2\pi}}\left[\tfrac{\Delta_i}{2\sigma}\exp\left(\tfrac{-\Delta_i^2}{8\sigma^2}\right)+\tfrac{\Delta_{i+1}}{2\sigma}\exp\left(\tfrac{-\Delta_{i+1}^2}{8\sigma^2}\right)\right] \right\} \\
			&= -H(X)-\tfrac12 + \sum_{i \in \calI} \left\{ \left(p_i\log \tfrac{1}{p_i}+\tfrac{p_i}{2}\right) \left[Q\left(\tfrac{\Delta_i}{2\sigma}\right)+Q\left(\tfrac{\Delta_{i+1}}{2\sigma}\right)\right] + \right.\\
			&\qquad \qquad \qquad \qquad \qquad \qquad  \qquad \qquad \qquad \qquad \qquad \left.\tfrac{p_i}{2\sqrt{2\pi}}\left[\tfrac{\Delta_i}{2\sigma}\exp\left(\tfrac{-\Delta_i^2}{8\sigma^2}\right)+\tfrac{\Delta_{i+1}}{2\sigma}\exp\left(\tfrac{-\Delta_{i+1}^2}{8\sigma^2}\right)\right] \right\}.
		\end{align*}
		Note that the $H(X) + \frac12$ from \eqref{eqn: h_x_y_ub} cancels. We now compute the second integral from \eqref{eqn: h_x_y_ub}, by using the upper bound $\log(1+z) \leq z$, as
		\begin{align*}
			& \quad \ \intR p_{i^*} \tfrac{1}{\sqrt{2\pi \sigma^2}}\exp\left(-\tfrac{{d^*}^2}{2\sigma^2}\right) \log \left[1 + \sum_{j \in \calI \setminus \{i^*\}} \frac{p_j}{p_{i^*}} \exp\left(\tfrac{{d^*}^2-d_j^2}{2\sigma^2}\right)\right] \d y\\
			&\leq \intR \sum_{j \in \calI \setminus \{i^*\}} p_j \tfrac{1}{\sqrt{2\pi \sigma^2}}\exp\left(-\tfrac{{d_j}^2}{2\sigma^2}\right) \d y\\
			&= \intR \sum_{j \in \calI} p_j \tfrac{1}{\sqrt{2\pi \sigma^2}}\exp\left(-\tfrac{{d_j}^2}{2\sigma^2}\right) \d y - \intR p_{i^*} \tfrac{1}{\sqrt{2\pi \sigma^2}}\exp\left(-\tfrac{{d^*}^2}{2\sigma^2}\right) \d y.
		\end{align*}
		The first term here is equal to 1, as interchanging the integral and summation over $j$ (by Fubini's theorem) gives 
		\begin{align*}
			\intR \sum_{j \in \calI} p_j \tfrac{1}{\sqrt{2\pi \sigma^2}}\exp\left(-\tfrac{{d_j}^2}{2\sigma^2}\right) \d y = \sum_{j \in \calI} p_j \intR \tfrac{1}{\sqrt{2\pi \sigma^2}}\exp\left(-\tfrac{{(y-x_j)}^2}{2\sigma^2}\right) \d y = \sum_{j \in \calI} p_j \bbE_{\tilde Y \sim \calN(x_j, \sigma^2)}[1] = 1.
		\end{align*}
		Meanwhile, the second term can be computed as
		\begin{align*}
			\intR p_{i^*} \tfrac{1}{\sqrt{2\pi \sigma^2}}\exp\left(-\tfrac{{d^*}^2}{2\sigma^2}\right) \d y &= \sum_{i \in \calI} p_i \int_{V_i} \tfrac{1}{\sqrt{2\pi \sigma^2}}\exp\left(-\tfrac{(y-x_i)^2}{2\sigma^2}\right) \d y\\
			&= \sum_{i \in \calI} p_i \left[Q\left(\tfrac{-\Delta_i}{2\sigma}\right)-Q\left(\tfrac{\Delta_{i+1}}{2\sigma}\right)\right]\\
			&= \sum_{i \in \calI} p_i \left[1 - Q\left(\tfrac{\Delta_i}{2\sigma}\right)-Q\left(\tfrac{\Delta_{i+1}}{2\sigma}\right)\right]\\
			&= 1 - \sum_{i \in \calI} p_i \left[Q\left(\tfrac{\Delta_i}{2\sigma}\right)+Q\left(\tfrac{\Delta_{i+1}}{2\sigma}\right)\right],
		\end{align*}
		or equivalently, the second integral from \eqref{eqn: h_x_y_ub} is upper bounded by $\sum_{i \in \calI} p_i \left[Q\left(\tfrac{\Delta_i}{2\sigma}\right)+Q\left(\tfrac{\Delta_{i+1}}{2\sigma}\right)\right]$. Putting all these together into \eqref{eqn: h_x_y_ub}, we have
		\begin{align*}
			H(X \mid Y) \leq \sum_{i \in \calI} \Big\{ \left(p_i\log \tfrac{1}{p_i}+\tfrac{3p_i}{2}\right) \left[Q\left(\tfrac{\Delta_i}{2\sigma}\right)+Q\left(\tfrac{\Delta_{i+1}}{2\sigma}\right)\right] &+ \frac{p_i}{2\sqrt{2\pi}}\Big[\tfrac{\Delta_i}{2\sigma}\exp\left(\tfrac{-\Delta_i^2}{8\sigma^2}\right)\\
			&\quad + \tfrac{\Delta_{i+1}}{2\sigma}\exp\left(\tfrac{-\Delta_{i+1}^2}{8\sigma^2}\right)\Big] \Big\}.
		\end{align*}
		We know that $Q(z)$ is decreasing for all $z$ and $z\exp(-\frac{z^2}{2})$ is decreasing for all $z > 1$. As $\dmin > 0$, for any $\sigma < \dmin(X)/2 \leq \Delta_i/2$ for all $i$, we have that each of the above $Q$ and $\exp$ terms is decreasing in its argument. Hence, we can further upper bound $H(X \mid Y)$ as
		\begin{align*}
			H(X \mid Y) &\leq Q\left(\tfrac{\dmin}{2\sigma}\right)(2H(X)+3) + \frac{1}{\sqrt{2\pi}}\tfrac{\dmin}{2\sigma}\exp\left(\tfrac{-\dmin^2}{8\sigma^2}\right)\\
			&\leq \exp\left(\tfrac{-\dmin^2}{8\sigma^2}\right) \left[H(X) + \tfrac32 + \tfrac1{\sqrt{2\pi}} \tfrac{\dmin}{2\sigma}\right],
		\end{align*}
		using the fact that $Q(z) \leq \frac{1}{2}\exp(-\frac{z^2}{2})$. This gives us the desired upper bound to $\lim_{\sigma \to 0}\sigma^2 \log H(X \mid Y)$, as
		\begin{align*}
			\lim_{\sigma \to 0} \sigma^2 \log H(X \mid Y) \leq \tfrac{-\dmin^2}{8} + \lim_{\sigma \to 0} \sigma^2 \log\left(H(X) + \tfrac32 + \tfrac1{\sqrt{2\pi}} \tfrac{\dmin}{2\sigma}\right) = \frac{-\dmin^2}{8},
		\end{align*}
		which completes the proof.
	\end{proof}

\end{document}
